% Chapter 1, typeset from lectures/L01: the README is the opener and the closing review, and
% each appendix is a section, in the same order and with the same subsections. The README's lecture
% plan is the plan for the live hour rather than the material, and is not typeset.

\chapter{Embedded Linux, and the Licence You Ship With It}
\label{c1:ch}

{\large\itshape What an embedded Linux system is made of, how it assembles itself at power-on, and
what the licences on the pieces oblige you to hand over.\par}

\begin{chapterbox}{In this chapter}
\begin{itemize}
  \item What is actually on the flash of a product that runs Linux, in four pieces, and what each
    one does at power-on.
  \item The boot chain from the on-chip ROM to the first userspace process, and the three steps
    QEMU's \code{-kernel} lets you skip.
  \item Cross-compilation: why \code{gcc} on your laptop is the wrong compiler, and what a sysroot
    is for.
  \item Buildroot and Yocto in one section, which is all they get in this course, and why.
  \item Copyleft against permissive, and the one question a licence actually answers.
  \item The kernel's syscall exception, \code{MODULE_LICENSE}, and a licence boundary the linker
    enforces.
  \item \code{make kernel}, \code{make qemu}, \code{make rootfs}, \code{make boot}, and a shell on a
    machine you just built.
\end{itemize}
\end{chapterbox}

This chapter is about what a product that runs Linux is made of, and what shipping it obliges you
to hand over. It takes the system apart into four pieces, the bootloader, the kernel, the root
filesystem and the toolchain that builds the other three, and follows the chain of programs that
assembles them between power-on and the first userspace process, including the three steps QEMU's
\code{-kernel} skips and what skipping them costs. It explains why the compiler on your laptop is
the wrong one and what a sysroot is for. Then it asks the one question a licence answers, what you
must give somebody along with a binary, works it through for the licences an embedded product
carries, and follows the two boundaries that matter for the rest of the book: the syscall exception,
which is why your application is yours, and \code{EXPORT_SYMBOL_GPL}, where the argument about
modules becomes a build error.

Read \secref{c1:app:a} first; it is the anatomy and the boot chain, and it ends with the machine
this course builds, piece by piece and with its sizes, in \secref{c1:app:a7}. That machine is
what the lab checks: what \code{make kernel}, \code{make qemu}, \code{make rootfs} and
\code{make boot} produce, rather than code you write. \Secref{c1:app:b} is the
licensing material, the one section in the book with no code in it and the one most likely to come
up in a design review. The exercises in \secref{c1:sec:exercises} end with the Cross-check: classify
the licences of ten kernel source files by hand, predict the whole directory from them, then run
your own audit script over it and reconcile the two, which will not agree. The lab is done when
\code{make test L=L01} reports \textbf{PASSED}, with all nine of the target's checks on the machine
you built, from the kernel version to the \code{qa-dev} node in its device tree.

% Section 1.1, from lectures/L01/appendix/a_the_four_pieces.md.

\section{The Four Pieces, and How They Assemble Themselves}
\label{c1:app:a}

This section is the anatomy of an embedded Linux system: what is actually stored on a product's
flash, what each piece does, and the order in which they hand control to one another between power
being applied and a prompt appearing. It ends with the machine this course runs on, and an honest
account of which parts of the chain that machine skips.

\subsection{What is on the flash}
\label{c1:app:a1}

An embedded Linux product contains four things. Three of them ship; the fourth never leaves your
build machine and determines everything about the other three.

\begin{booktable}{@{}l>{\hsize=0.49\hsize}L>{\hsize=1.51\hsize}L@{}}
  \thead{Piece} & \thead{Ships on the device} & \thead{What it does} \\
  \midrule
  Bootloader & Yes & Brings up DRAM, finds the kernel, hands it a device tree \\
  Kernel & Yes & Drives the hardware and provides the system call interface \\
  Root filesystem & Yes & Everything that runs as a process, starting with PID 1 \\
  Toolchain & No & Compiles the other three; fixes the ABI they share \\
\end{booktable}

\noindent Almost every question in the rest of this course is a question about which of the four a
thing belongs to. ``Why does my driver not load'' is a kernel question. ``Why is my binary 40 MB''
is a toolchain question. ``Why does nothing happen after \code{Run /init as init process}'' is a
root filesystem question, and the answer is nearly always that \file{/init} is not there, is not
executable, or is dynamically linked against a library that is not there either.

The fourth piece is the one people forget. A toolchain is not just a compiler; it is a compiler, a C
library, a linker, and a set of headers, and the choice of C library in particular is baked into
every binary you ship. You cannot mix a glibc binary and a musl root filesystem, and the failure
when you try is a loader error at run time and not a link error at build time.

\subsection{The boot chain}
\label{c1:app:a2}

Power is applied. What happens next is a chain of five programs, each one loading the next, and each
one existing because the one before it did not have enough memory to do the next one's job.

\begin{console}
  power on
     |
     v
  [1] Boot ROM          on-chip, mask-programmed, a few kilobytes
     |                  reads a strap pin, loads the next stage into on-chip SRAM
     v
  [2] SPL / first stage  ~64 KB, runs from SRAM
     |                  its whole purpose is to initialise the DRAM controller
     v
  [3] Bootloader         ~1 MB, runs from DRAM (U-Boot, or an EFI loader)
     |                  loads the kernel and the DTB into DRAM, then jumps
     v
  [4] Kernel             decompresses itself, drives the hardware
     |                  mounts a root filesystem, then hands over
     v
  [5] /init              PID 1, the first userspace process
                        starts everything else and never exits
\end{console}

\noindent Two things about this chain are worth stating explicitly because they are the parts people
get wrong.

\textbf{The SPL exists because of a chicken-and-egg problem.} U-Boot is around a megabyte and needs
DRAM to run in. DRAM does not work until its controller has been configured, with a long sequence of
board-specific timing values. That configuration code has to run from somewhere, and the only memory
available at that point is the few tens of kilobytes of SRAM on the SoC itself. So the build
produces a second, tiny bootloader whose entire job is to set up DRAM and then load the real one. If
you have ever wondered why a board has both \code{MLO} and \file{u-boot.img}, or both \code{SPL} and
\code{u-boot}, that is why.

\textbf{The boot ROM is the only piece you cannot replace.} It is mask-programmed into the silicon.
This is what makes bricking recoverable on some boards and terminal on others: if the ROM will fall
back to loading over USB or UART when the flash is empty, you can always recover; if it will not, a
bad write to the bootloader partition turns the board into a coaster.

\lecturefigure{lectures/L01/appendix/images/boot_chain.png}
  {The boot chain drawn as five boxes, each loading the next: ROM code on the die, SPL in SRAM, U-Boot
   bringing up DRAM, the kernel Image and DTB, and init as PID 1. The boxes are coloured by which of
   the four pieces supplies them, with the toolchain named as the piece that never ships.}
  {fig:boot_chain}

\subsection{What the kernel is handed, and by whom}
\label{c1:app:a3}

Step 3 hands step 4 two things: the kernel image, and a \textbf{device tree blob}.

The device tree matters enough to have a chapter of its own (Chapter~\ref{c10:ch}), but the reason
it exists belongs here. An x86 PC can discover its own hardware: PCI is enumerable, ACPI describes
the rest, and a single kernel image boots on any machine. An embedded SoC can do none of that. There
is no bus to walk that will tell you a UART lives at \code{0x09000000} and raises interrupt 33.
Somebody has to say so, and the device tree is a data file that says so, kept separate from the
kernel so that one kernel image can boot many boards.

So the kernel receives, in a register, a pointer to a blob that describes the machine. It parses it,
and every driver that later probes gets its addresses from it. That is the whole idea, and
Chapter~\ref{c10:ch} is about the mechanism.

\subsection{What this course skips, and what that costs}
\label{c1:app:a4}

\code{make boot} runs this:

\begin{shell}
qemu-system-aarch64 -machine virt -kernel Image -initrd rootfs.cpio.gz -append "..."
\end{shell}

\noindent \code{-kernel} skips steps 1, 2 and 3 entirely. QEMU writes the kernel image into the
guest's memory, generates a device tree describing the machine it is emulating, puts a pointer to it
in the right register, and jumps to the kernel's entry point. It is doing U-Boot's job, without
being U-Boot.

\textbf{What that costs you is real and worth naming.} You will not learn here how to configure a
DRAM controller, how to write a U-Boot environment, how to recover a board whose bootloader you
overwrote, or what the serial output of a failing SPL looks like. Those are important skills and
this course does not teach them.

\textbf{What it buys is that everything after step 3 is identical.} The kernel does not know or care
whether U-Boot or QEMU put it in memory; it reads the same device tree and probes the same drivers
either way. Every chapter from Chapter~\ref{c3:ch} onwards would be word-for-word the same on a real
board.

\subsection{The toolchain, and why your compiler is the wrong one}
\label{c1:app:a5}

The target runs 64-bit ARM. Your build machine almost certainly does not. So the compiler that
produces target binaries is not \code{gcc} but \code{aarch64-linux-gnu-gcc}, and that prefix is a
\textbf{triplet}, or in this case something closer to a quadruplet:

\begin{console}
aarch64  -  linux  -  gnu
   |          |        |
   |          |        +--  C library and ABI: gnu (glibc). Others: musl, uclibc.
   |          +-----------  operating system: linux. Others: none, elf, darwin.
   +----------------------  architecture: aarch64. Others: arm, x86_64, riscv64.
\end{console}

\noindent The single most useful demonstration of what that means is one command:

\begin{shell}
$ gcc -o hello-host hello.c && file hello-host
hello-host: ELF 64-bit LSB pie executable, x86-64, ...

$ aarch64-linux-gnu-gcc -o hello-target hello.c && file hello-target
hello-target: ELF 64-bit LSB pie executable, ARM aarch64, ...
\end{shell}

\noindent Same source, same command line, two files that share nothing. The second will not run on
your laptop and the first will not run on the target.

\textbf{The sysroot is the other half.} A cross-compiler must not use your machine's libraries,
because those describe your machine. It uses a \emph{sysroot}: a directory that mirrors the target's
filesystem, holding the target's headers and the target's libraries. A toolchain built by Buildroot
or Yocto names that directory when asked. Debian's, the one in the course's container, answers
\code{/}:

\begin{shell}
$ aarch64-linux-gnu-gcc -print-sysroot
/
\end{shell}

\noindent That is because Debian keeps the target's files beside your machine's rather than in a
tree of their own: the target's C library and headers in \file{/usr/aarch64-linux-gnu}, other arm64
libraries in \file{/usr/lib/aarch64-linux-gnu}, and headers that are the same on every architecture
in \file{/usr/include}, which this compiler searches too. What it never searches is
\file{/usr/lib/x86_64-linux-gnu}, where your machine's own libraries are.

Two consequences follow, and both of them bite people. A library your program needs must be
installed \emph{for the target}, not for your machine; installing \code{libfoo-dev} with apt
installs your machine's library, which the cross linker cannot use. And the target's glibc version
sets the oldest target root filesystem your binary will run on, because glibc guarantees that a
binary runs on a newer glibc than the one it was built against, and not on an older one.

\textbf{Kernel code has no sysroot and no C library at all.} It is compiled \code{-nostdinc},
against the kernel's own headers. There is no \code{printf}, no \code{malloc}, and no \code{double}.
Chapter~\ref{c4:ch} covers what that environment is actually like; the point here is that the kernel
does not use the C library the rest of the toolchain is built around, which is why ``the toolchain''
and ``the kernel'' are two of the four pieces and not one.

\subsection{Buildroot and Yocto, in one section}
\label{c1:app:a6}

The three shipping pieces have to be built, configured and packaged together, reproducibly, by
somebody who is not you and possibly on a build server. That is a real problem and there are two
mainstream answers to it.

\begin{booktable}{@{}l>{\hsize=0.84\hsize}L>{\hsize=1.16\hsize}L@{}}
   & \thead{Buildroot} & \thead{Yocto / OpenEmbedded} \\
  \midrule
  Model & One \code{make}, one config, one image & Layers of recipes, one per package \\
  Output & A root filesystem image & A root filesystem image, and a package feed \\
  Rebuild cost & Often a full rebuild & Incremental, per recipe \\
  Learning cost & An afternoon & Weeks \\
  Fits & Small, fixed products & Products with variants, updates, or a team \\
\end{booktable}

\noindent The honest summary is that Buildroot is a build system and Yocto is a distribution
builder, and teams pick the wrong one about half the time. A product with one variant and no field
updates does not need Yocto's machinery. A product with four hardware variants, two customers and an
OTA update channel will outgrow Buildroot's single-config model.

\textbf{This course uses neither.} Its root filesystem is one statically linked BusyBox and a shell
script, assembled by \file{ci/rootfs.sh} in about a hundred lines. That is not what a product does, and
it is not held up as an example of what to do. It is the smallest thing that boots, which makes it
the right thing to learn on, and it means every file on the target is one you can account for.

\subsection{The machine this course builds}
\label{c1:app:a7}

Running \code{make kernel}, \code{make qemu}, \code{make rootfs} and \code{make boot} produces this:

\begin{booktable}{@{}lLl@{}}
  \thead{Piece} & \thead{What it is here} & \thead{Size} \\
  \midrule
  Bootloader & None. QEMU's \code{-kernel} does the job & 0 \\
  Kernel & \code{Image}, Linux 6.12.30, cross-built for arm64 & \textasciitilde{}24 MB \\
  Root filesystem & \code{rootfs.cpio.gz}: BusyBox, static, plus three modules & \textasciitilde{}1.3 MB \\
  Toolchain & \code{aarch64-linux-gnu-gcc} 12, from the container image & n/a \\
\end{booktable}

\noindent The two sizes are worth staring at for a moment, because the ratio is the wrong way round
from what most people expect. \textbf{The kernel is roughly twenty times the size of the entire
userland.} That is not because the kernel is bloated in some abstract sense; it is because the
configuration it was built from still enables a great deal that this machine cannot use. Finding out
how much, and taking it out, is Chapter~\ref{c3:ch}'s lab.

The exercises in \S\ref{c1:sec:exercises} ask you to produce both of those numbers yourself rather
than take them from this table, and the Cross-check asks you to predict one of them before you
measure it.

\subsection{Where to look on a running system}
\label{c1:app:a8}

Once \code{make boot} gives you a prompt, four files tell you which of the four pieces you are
looking at. Chapter~\ref{c2:ch} is about this in earnest; these are the ones worth knowing from the
start.

\begin{console}
/proc/version      the kernel: version, compiler, and when it was built
/proc/cmdline      what the bootloader (here, QEMU) told the kernel at handover
/proc/iomem        the physical address map, as the kernel understands it
/sys/firmware/devicetree/base
                   the device tree, unpacked into a directory per node
\end{console}

\noindent That last one is worth a minute now even though Chapter~\ref{c10:ch} is where it is
explained. The device tree the kernel was handed is exposed as a directory tree, one directory per
node and one file per property. The device this course is built around appears there:

\begin{console}
/sys/firmware/devicetree/base/platform-bus@c000000/qa-dev@0/
    compatible      "qacademy,qa-dev-1.0"
    reg             00 00 00 00 00 00 10 00
    interrupts      00 00 00 00 00 00 00 70 00 00 00 04
\end{console}

\noindent Nobody typed those numbers. QEMU chose where to put the device, wrote the node, and handed
the result to the kernel, which is exactly what a bootloader does on a real board. Reading them back
and working out what physical address they describe is Chapter~\ref{c6:ch}'s Cross-check.

% Section 1.2, from lectures/L01/appendix/b_licences.md.

\section{Open Source Licences, and What You Have to Publish}
\label{c1:app:b}

This section is about one question. It is the only question a licence actually answers, and almost
every argument about open source licensing is an argument that has lost sight of it:

\begin{quote}
\textbf{If I give somebody this binary, what am I obliged to give them along with it?}
\end{quote}

Not ``is this free software''. Not ``can I use this commercially''; you almost always can. The
question is what leaves your building when the product does.

\textbf{This is not legal advice, and the author is not a lawyer.} What it is is the engineering
model you need to have in your head so that you know which decisions need a lawyer and which do not.
Getting this wrong is expensive, and the expensive cases are almost always cases where nobody asked
the question at all until a customer did.

\subsection{Two families}
\label{c1:app:b1}

Every licence you will meet on an embedded product falls into one of two families, distinguished by
what happens when you combine the code with your own.

\textbf{Permissive} licences let you do essentially anything, including shipping binaries with no
source, provided you carry the copyright notice along. MIT, BSD-2-Clause, BSD-3-Clause and
Apache-2.0 are the ones you will actually see.

\textbf{Copyleft} licences let you do essentially anything, provided that whoever receives the
binary can also receive the source, under the same licence. GPL-2.0, GPL-3.0, LGPL-2.1 and MPL-2.0
are the ones you will see.

\begin{booktable}{@{}l>{\hsize=0.44\hsize}L>{\hsize=1.56\hsize}L@{}}
  \thead{Licence} & \thead{Family} & \thead{If you ship a binary containing it, you must publish} \\
  \midrule
  MIT & Permissive & Nothing. Keep the notice in your documentation \\
  BSD-3-Clause & Permissive & Nothing, plus you may not use the authors' names to endorse \\
  Apache-2.0 & Permissive & Nothing, plus a NOTICE file, plus a patent grant you give \\
  MPL-2.0 & Weak copyleft (file) & The source of the MPL files you changed \\
  LGPL-2.1 & Weak copyleft (library) & The library's source, and the means to relink your binary against a modified copy \\
  GPL-2.0 & Strong copyleft & The complete corresponding source of the whole work \\
  GPL-3.0 & Strong copyleft & The above, plus installation information for consumer devices \\
\end{booktable}

\noindent The two rows worth reading twice are LGPL and GPL-3.0, because both carry an obligation
that is not about source code at all, and both are where embedded products get caught.

\subsection{The word that does the work: derivative}
\label{c1:app:b2}

Copyleft only reaches code that is part of the same work. So everything depends on where one work
ends and the next begins, and that boundary is not defined by the licence in terms an engineer can
apply mechanically. What exists instead is a set of positions, some of which are near-universally
accepted and some of which are contested.

Broadly accepted:
\begin{itemize}
  \item \textbf{Static linking creates one work.} If your proprietary code is linked into the same
    binary as GPL code, the result is a derivative work of both.
  \item \textbf{Separate processes communicating over a pipe or a socket do not.} Your closed daemon
    talking to a GPL program over a socket is two works.
  \item \textbf{Mere aggregation does not.} Two unrelated programs on the same filesystem image are
    not one work because they share a partition.
\end{itemize}
Contested, and where you need a lawyer rather than an opinion:
\begin{itemize}
  \item \textbf{Dynamic linking.} The Free Software Foundation holds that linking against a GPL
    library creates a derivative work even at run time. Others disagree. This is why the LGPL
    exists: it settles the question for libraries that want to be linkable by proprietary code.
  \item \textbf{Kernel modules.} See \S\ref{c1:app:b4}, which is the case this course actually cares
    about.
\end{itemize}

\subsection{The syscall exception, and why your application is yours}
\label{c1:app:b3}

The Linux kernel is GPL-2.0-only. Every program on a Linux system makes system calls into it. If
that made every program a derivative work of the kernel, no proprietary software could run on Linux
at all, and Linux would have no commercial users.

Linus Torvalds added a note to the top of the kernel's \file{COPYING} file that says so explicitly.
The note now lives in \file{LICENSES/exceptions/Linux-syscall-note}, which \file{COPYING} names as
the kernel's explicit syscall exception:

\begin{quote}
NOTE! This copyright does \emph{not} cover user programs that use kernel services by normal system
calls --- this is merely considered normal use of the kernel, and does \emph{not} fall under the
heading of ``derived work''.
\end{quote}

That paragraph is the legal foundation of embedded Linux as a commercial platform, and it is worth
knowing that it is a clarification granted by the copyright holder rather than something that falls
out of the GPL itself. \textbf{The boundary it draws is the system call boundary.} Your application
talks to the kernel through \code{read}, \code{write}, \code{ioctl} and the rest, and that is normal
use.

The practical consequence is the one your colleague gets wrong. ``We ship Linux, so all our code
must be open source'' is false: your application, running as its own process and making system
calls, is yours, and you may ship it as a binary with no source at all. What \emph{is} true is the
narrower claim, that the kernel you ship, and any modifications you made to it, must be published.

\subsection{Kernel modules, and the licence the linker enforces}
\label{c1:app:b4}

A kernel module is not a separate process. It is compiled against the kernel's own headers, linked
into the kernel's address space, and calls kernel functions directly. Under every version of the
derivative-work argument, that is a much harder position to defend than an application making system
calls, and the widely held view among kernel developers is that a module is a derivative work of the
kernel.

The kernel does not leave this as an argument. It enforces a version of it mechanically, and this is
the part that matters to you as a driver author.

\textbf{Every module declares its licence in its own binary:}

\begin{ccode}
MODULE_LICENSE("GPL");
\end{ccode}

\noindent The string is not documentation. It is read by the module loader, and the set of accepted
values is fixed: \code{"GPL"}, \code{"GPL v2"}, \code{"GPL and additional rights"}, \code{"Dual
BSD/GPL"}, \code{"Dual MIT/GPL"}, \code{"Dual MPL/GPL"}, and \code{"Proprietary"}. Anything else
counts as proprietary. A module with no \code{MODULE_LICENSE} at all is not built: \code{modpost}
fails with \code{missing MODULE_LICENSE()}.

\textbf{Two things follow when a module is not GPL-compatible.}

First, the kernel is \textbf{tainted}. \file{/proc/sys/kernel/tainted} gains a bit, every subsequent
oops carries a taint flag, and kernel developers will decline to debug it. That is a social
mechanism rather than a legal one, but it is a real cost: a bug report from a tainted kernel is a
bug report nobody upstream will look at.

Second, and concretely, \textbf{a large part of the kernel's internal API becomes unavailable.}
Symbols are exported in one of two ways:

\begin{ccode}
EXPORT_SYMBOL(foo);      /* any module may use this */
EXPORT_SYMBOL_GPL(bar);  /* only GPL-compatible modules may use this */
\end{ccode}

\noindent A module declaring a proprietary licence and calling \code{bar()} is refused, and it is
refused \textbf{in two different places depending on what the build could see.}

Usually the build fails, at the \code{modpost} step that runs after linking each module:

\begin{console}
ERROR: modpost: GPL-incompatible module mine.ko uses GPL-only symbol 'bar'
\end{console}

\noindent That is the good case: the message names the licence, the module and the symbol, and no
\code{.ko} is produced. It happens whenever \code{modpost} can see that \code{bar} is GPL-only,
which is nearly always, because the built-in kernel's exports are listed in its
\file{Module.symvers}.

When \code{modpost} could not have known, the module loader refuses instead, and its message is much
worse. It does not report a licence problem; it \textbf{hides GPL-only symbols from a proprietary
module's symbol lookup entirely}, so what you see is:

\begin{console}
mine: Unknown symbol bar (err -2)
\end{console}

\noindent which is indistinguishable from a typo or a missing dependency. That path is what you hit
after a kernel upgrade in which a symbol you used changed from \code{EXPORT_SYMBOL} to
\code{EXPORT_SYMBOL_GPL}; your module was built when it was legal and is loaded when it is not.
Chapter~\ref{c4:ch} causes both on purpose.

\textbf{Which symbols are \code{_GPL} is a live decision, not a historical accident.} Maintainers
choose, and over the years the boundary has moved steadily toward \code{_GPL} for anything that
looks like a core internal interface. If your business model depends on a proprietary module, it
depends on a set of exported symbols that other people control and may narrow in the next release.

\subsection{What ``complete corresponding source'' means in practice}
\label{c1:app:b5}

GPL-2.0 section 3 requires you to give the recipient ``the complete corresponding machine-readable
source code'', defined as ``all the source code for all modules it contains, plus any associated
interface definition files, plus the scripts used to control compilation and installation of the
executable''.

Read the last clause again, because it is the one people fail. \textbf{The build scripts are part of
the obligation.} A tarball of kernel sources with no \code{.config} in it does not satisfy this: the
recipient cannot produce your binary from it. Neither does a source tree that only builds with an
internal toolchain nobody outside your company has.

The practical checklist for a product that ships a GPL kernel is:
\begin{itemize}
  \item The kernel source, at the exact version you shipped, including every patch you applied.
  \item The \code{.config} you built it with.
  \item The device tree sources.
  \item Enough instructions to build it: which toolchain, which commands.
  \item The same for U-Boot, BusyBox, and every other GPL component in the image.
\end{itemize}
You may either ship this alongside the product, or include a \textbf{written offer} valid for three
years to supply it on request. The written offer is what most products use, and it is the thing that
has to survive in your build system for three years after the last unit ships. That is a
release-engineering requirement that comes directly out of a licence, and it is the single most
common thing to discover too late.

\textbf{GPL-3.0 adds installation information.} For a consumer device, you must also provide
whatever is needed to actually install a modified version and have it run: keys, signing tools,
flashing instructions. This is why a great many embedded products deliberately stay on GPL-2.0-only
components, and why the kernel's own GPL-2.0-only status without the ``or later'' clause is a
decision rather than an oversight.

\subsection{SPDX, and how you find out what you have}
\label{c1:app:b6}

You cannot answer the question in the opening of this section without an inventory. Reading every
file is not possible; almost every project now says so in a machine-readable way instead.

An \textbf{SPDX identifier} is a single comment line at the top of a file:

\begin{ccode}
// SPDX-License-Identifier: GPL-2.0-only
\end{ccode}

\noindent The kernel adopted these in 2017 and nearly every file now carries one. They are exact,
they are greppable, and a file that has one needs no further interpretation:

\begin{shell}
grep -rh 'SPDX-License-Identifier:' drivers/ | sort | uniq -c | sort -rn
\end{shell}

\noindent Three cautions, and they are the whole of what makes this a skill rather than a
\code{grep}.
\begin{itemize}
  \item \textbf{A file with no SPDX tag is not unlicensed.} It is a file whose licence has to be
    found some other way: a \code{LICENSE} file in its directory, a header comment, or the project's
    top-level terms. These are the interesting files in any audit, and there are always some.
  \item \textbf{\code{GPL-2.0-only} and \code{GPL-2.0-or-later} are different licences.} The second
    lets a downstream recipient move to GPL-3.0 and its installation-information obligation. The
    kernel is \code{GPL-2.0-only}. Some of the components around it are not.
  \item \textbf{A build system is not an inventory.} What matters is what ended up in the image, and
    a package can be present in the source tree and absent from the shipped filesystem, or the
    reverse.
\end{itemize}
Exercise~\ref{c1:ex:6} is to build that inventory for a real tree and say what you would have to
publish. The Cross-check, Exercise~\ref{c1:ex:8}, is to do a subdirectory by hand first, and find
out where your script and your eyes disagree.

\subsection{A worked example}
\label{c1:app:b7}

A product ships:

\begin{narrowtable}{@{}lll@{}}
  \thead{Component} & \thead{Licence} & \thead{Modified} \\
  \midrule
  U-Boot & GPL-2.0-or-later & Yes, board port \\
  Linux kernel & GPL-2.0-only & Yes, one driver \\
  Your out-of-tree driver & Yours to choose & n/a \\
  BusyBox & GPL-2.0-only & No \\
  An MQTT client library & Apache-2.0 & No \\
  Your control application & Proprietary & n/a \\
\end{narrowtable}

\noindent What must you publish?
\begin{itemize}
  \item \textbf{U-Boot}, complete corresponding source including your board port and your config.
  \item \textbf{The kernel}, complete corresponding source including your driver, your
    \code{.config} and your device tree. Note that the driver is in the kernel tree and is therefore
    part of the work.
  \item \textbf{BusyBox}, source at the version shipped. Unmodified does not mean exempt; the
    obligation is to the recipient of the binary, not conditional on your having changed it.
  \item \textbf{The MQTT library}: nothing, but carry its NOTICE file and its licence text in your
    documentation.
  \item \textbf{Your control application}: nothing. It is a separate process making system calls,
    and \S\ref{c1:app:b3} is why.
\end{itemize}
And the one that is a decision rather than an obligation: \textbf{your out-of-tree driver}. If you
keep it out of tree and declare it proprietary, you take on the taint, you lose every
\code{EXPORT_SYMBOL_GPL} symbol, and you own the maintenance of it against a kernel API that has no
stable interface and changes every release. Most companies that try this eventually stop, and the
reason is usually the maintenance cost rather than the licence.

% Section 1.3, from lectures/L01/README.md: what you should be able to do afterwards, the
% questions to test yourself, and the next lecture.

\section{Review}
\label{c1:sec:review}

\needspace{6\baselineskip}
\subsection*{What you should be able to do now}
After this chapter, you should be able to:
\begin{itemize}
  \item Name the four pieces of an embedded Linux system, and say for any given file which one it
    belongs to.
  \item Walk the boot chain from power-on to the first userspace process, and say what each step
    does that the one before it could not.
  \item Say what \code{-kernel} skips, and therefore what this course does not teach you.
  \item Explain what a toolchain triplet names, and why a sysroot is not the same thing as the
    host's \file{/usr/include}.
  \item Answer the licence question for a given component: what must be published, to whom, and
    when.
  \item Say why an application that only makes system calls is not a derivative work of the kernel,
    and why a kernel module usually is.
  \item Predict whether a module will load, given its \code{MODULE_LICENSE} and the symbols it uses.
\end{itemize}

\needspace{6\baselineskip}
\subsection*{Questions to test yourself}
\begin{enumerate}
  \item A colleague says ``we ship Linux, so all our code has to be open source''. What is wrong
    with that sentence, and what is the part of it that is right?
  \item Your product has a bootloader, a kernel, a BusyBox root filesystem and one proprietary
    daemon you wrote. A customer asks for the source. What do you have to give them?
  \item Why does the boot chain have an SPL at all, given that U-Boot could do the same job?
  \item What does \code{aarch64-linux-gnu} tell you that \code{aarch64} alone does not?
  \item A driver is licensed MIT and calls a function exported with \code{EXPORT_SYMBOL_GPL}. What
    happens, and at which point: compile, link, load, or run?
  \item Why is the kernel image 24 MB when the machine it boots has one serial port and one timer?
\end{enumerate}

\needspace{6\baselineskip}
\subsection*{Where to look}
\begin{itemize}
  \item The build targets are documented in the repository's root \file{README.md}, and
    \code{make help} lists them all.
\end{itemize}

\needspace{6\baselineskip}
\subsection*{Looking ahead}
Chapter~\ref{c2:ch} is about:
\begin{itemize}
  \item The command line, taken as the set of questions you are able to ask a running kernel.
  \item \file{/proc} and \file{/sys}: not directories, and not files either.
  \item What a two-megabyte userland gives up, measured against the one on your laptop.
  \item Why PID 1 is different from every other process, and what happens if it exits.
\end{itemize}
The machine you built in this chapter is the machine the other eleven chapters run on.
Chapter~\ref{c2:ch} is about learning to ask it what it is doing.

% Section 1.4, from lectures/L01/appendix/c_exercises.md.

\section{Exercises}
\label{c1:sec:exercises}\label{c1:app:c}

\begin{aside}
Eight, ending with the Cross-check. Do them in order; the last one needs the script from
Exercise~\ref{c1:ex:6}.

Several of these have a plausible wrong answer that is worth walking into before you find it. Where
an exercise can be checked mechanically, a \textbf{Check yourself} line says how.

Everything here runs against the trees \code{make kernel} and \code{make rootfs} already downloaded,
under \file{build/}. Nothing needs the network.

\textbf{Solutions.} The exercises that are not Code are answered in \secref{sol:sec:c1}. The Code
exercises have no solution: each is checked by running it, as its Check yourself line says.
\end{aside}

% ------------------------------------------------------------------------------------------
\exercise[fillaccent2]{The four pieces}{Recall}
\label{c1:ex:1}

Without looking at \S\ref{c1:app:a}, name the four pieces of an embedded Linux system and say in one
sentence what each does.

Then, for each of the following, say which piece it belongs to:

\xpart{a} \file{/lib/modules/6.12.30/kernel/lib/kunit/kunit.ko}

\xpart{b} \file{u-boot.img}

\xpart{c} \code{aarch64-linux-gnu-gcc}

\xpart{d} \file{/init}

\xpart{e} The \code{.dtb} handed to the kernel at boot

\xpart{f} \file{/bin/busybox}

% ------------------------------------------------------------------------------------------
\exercise[fillaccent2]{Why there are two bootloaders}{Recall}
\label{c1:ex:2}

A colleague looking at a board's flash layout asks why there is both an \code{SPL} and a
\file{u-boot.img}, and suggests deleting the first to save space.

\xpart{a} Explain in two sentences why the SPL exists.

\xpart{b} What specifically would fail if you deleted it, and at what point in the boot?

\xpart{c} Name one thing about a board that makes bricking it recoverable, and one that does not.

% ------------------------------------------------------------------------------------------
\exercise[fillaccent3]{What the toolchain triplet promises}{Hand calculation}
\label{c1:ex:3}

\xpart{a} Break \code{aarch64-linux-gnu} into its parts and say what each names.

\xpart{b} You are handed a binary and told only that it was built with \code{arm-none-eabi-gcc}.
What can you say about the system it is intended to run on, and in particular, can it run on Linux?
Justify it from the triplet alone.

\xpart{c} A program builds against \code{aarch64-linux-gnu-gcc} on a machine whose sysroot has glibc
2.36. It is deployed to a target whose root filesystem has glibc 2.31. Predict what happens, and say
whether the failure is at build time or at run time.

\xpart{d} A colleague fixes a ``missing header'' error in a cross build by running \code{sudo apt
install libfoo-dev}. Explain why this does not work, and what would.

% ------------------------------------------------------------------------------------------
\exercise[fillaccent4]{What leaves the building}{Design}
\label{c1:ex:4}

You ship a product containing exactly these things:

\begin{booktable}{@{}>{\hsize=1.06\hsize}Ll>{\hsize=0.94\hsize}L@{}}
  \thead{Component} & \thead{Licence} & \thead{Modified by you} \\
  \midrule
  U-Boot & GPL-2.0-or-later & Yes, a board port \\
  Linux kernel & GPL-2.0-only & Yes, one driver in-tree \\
  BusyBox & GPL-2.0-only & No \\
  \code{libcurl} & MIT-like (curl) & No \\
  A logging daemon you wrote & Proprietary & n/a \\
  A kernel module you wrote & You choose & n/a \\
\end{booktable}

\noindent A customer sends a written request for the source code.

\xpart{a} Go through the table and say, for each row, what you must publish.

\xpart{b} BusyBox is unmodified. Does that change the answer? Say why.

\xpart{c} Your logging daemon links dynamically against \code{libcurl} and talks to your kernel
module through \code{ioctl}. Does either of those relationships oblige you to publish the daemon?

\xpart{d} What must you keep, and for how long, in order to be able to answer this request at all?

% ------------------------------------------------------------------------------------------
\exercise[fillaccent4]{Where the boundary falls}{Design}
\label{c1:ex:5}

For each pair below, say whether the two things are one work or two, and name the argument that
decides it. Where the answer is genuinely contested, say so rather than picking.

\xpart{a} A proprietary application and the kernel, communicating by system calls.

\xpart{b} A proprietary program statically linked against a GPL-2.0 library.

\xpart{c} A proprietary program dynamically linked against an LGPL-2.1 library.

\xpart{d} A proprietary program dynamically linked against a GPL-2.0 library.

\xpart{e} A proprietary kernel module and the kernel.

\xpart{f} A GPL program and a proprietary program on the same filesystem image, which never
communicate.

% ------------------------------------------------------------------------------------------
\exercise{An SPDX inventory}{Code}
\label{c1:ex:6}

Write a script, \file{lectures/L01/lab/spdx-audit.sh}, that takes a directory and reports what
licences the files under it are under.

It should:
\begin{enumerate}[label=\textbf{\alph*)}]
  \item Walk every \code{*.c} and \code{*.h} file under the directory it is given.
  \item For each, extract the SPDX identifier if there is one. The tag looks like
    \code{SPDX-License-Identifier: GPL-2.0-only} and appears in a comment in roughly the first few
    lines.
  \item Print a count per identifier, most common first.
  \item Print, separately and by name, every file that has no identifier at all. These are the
    interesting ones and they must not be silently folded into a total.
  \item Print the total number of files examined, so that the counts can be checked against it.
\end{enumerate}
Two requirements that are easy to miss and are the point of the exercise:
\begin{itemize}
  \item \textbf{Normalise the deprecated spellings.} SPDX renamed several identifiers;
    \code{GPL-2.0} is the old spelling of \code{GPL-2.0-only}, and \code{GPL-2.0+} is the old
    spelling of \code{GPL-2.0-or-later}. Both spellings are in the kernel tree, in quantity. A
    report that lists them separately is reporting two licences where there is one. Your script must
    report the current spelling and say how many of each it merged.
  \item \textbf{Do not stop at the first match in the file.} Some files carry more than one SPDX
    tag. Take the first one, but know that you have made that choice.
\end{itemize}

\checkyourself run it on \file{build/linux-6.12.30/drivers/rtc}, which has 184 \code{.c} and
\code{.h} files at its top level. Your total must be 184.

% ------------------------------------------------------------------------------------------
\exercise{Measure the machine you built}{Code}
\label{c1:ex:7}

Answer each of these with a command, run on the target after \code{make boot}, or on the host
against \file{build/}. Write down the command as well as the number.

\xpart{a} What kernel version is running, what compiler built it, and when?

\xpart{b} What did the bootloader pass to the kernel on its command line?

\xpart{c} How large is the kernel image, and how large is the compressed initramfs?

\xpart{d} How many modules did the kernel build, and how many are installed on the target?

\xpart{e} How much memory does the target have, and how much of it is free once it has booted?

\xpart{f} What is the physical address of the \code{qa-dev} registers, according to
\file{/proc/iomem}?

For \textbf{f)}, note that nothing has claimed the region yet, because no driver for the device
exists until Chapter~\ref{c6:ch}. Say what you can see and what you cannot.

% ------------------------------------------------------------------------------------------
\exercise[fillaccent]{An audit by hand against an audit by script}{Cross-check}
\label{c1:ex:8}

This is the exercise the other seven exist for. You will compute a set of numbers by hand, run your
own code on the same problem, and reconcile the two. They will not agree, and the disagreement is
the lesson.

\xpart{a} By hand, with no scripting, classify the licences of these ten files in
\file{build/linux-6.12.30/drivers/rtc}:

\begin{console}
rtc-ds1307.c   rtc-pl031.c    rtc-starfire.c   rtc-rs5c313.c   rtc-cmos.c
rtc-m41t80.c   rtc-pcf8563.c  rtc-abx80x.c     rtc-max77686.c  rtc-s35390a.c
\end{console}

\noindent For each, write down the identifier and where you found it. Do not run your script yet.

\xpart{b} Now predict, from that sample of ten, what the breakdown over the whole of
\file{drivers/rtc} will be. Write down a number of distinct licences and a rough proportion for
each.

\xpart{c} Run your script from Exercise~\ref{c1:ex:6} over the whole directory. Record the output.

\xpart{d} Reconcile. Specifically, answer these:
\begin{itemize}
  \item How many \emph{distinct SPDX strings} did your script find, and how many \emph{distinct
    licences} is that actually? What is the relationship between the two numbers?
  \item Your script reports some number of untagged files. Find each one and determine its licence
    by another route. Say what that route was, and whether it was the same route for both.
  \item Did your ten-file sample predict the whole? If your proportions were off, say which
    direction and why a sample drawn from files you recognised the names of would be biased.
\end{itemize}

\xpart{e} A colleague runs your script, sees no \code{Proprietary} in the output, and concludes that
the kernel tree contains nothing proprietary. Say what is wrong with that inference. There are at
least two separate problems with it.

